\chapter{Vector Calculus and Maxwell Notation}
\label{app:vector}

This appendix collects the vector-calculus identities, integration
formulas, and notational conventions used throughout the book. All
results are stated for three-dimensional space with Cartesian, or
where noted, curvilinear coordinates.


%% ================================================================
\section{Curl identities}
\label{sec:curl-identities}

The following identities are used repeatedly in the derivation of
Maxwell inner products and self-adjointness conditions.

\paragraph{Double curl.}
\begin{equation}
 \curl(\curl\FF) = \grad(\divg\FF) - \nabla^2\FF.
 \label{eq:double-curl}
\end{equation}

\paragraph{Curl of a product.}
\begin{equation}
 \curl(f\,\FF) = f\,\curl\FF + (\grad f)\times\FF.
 \label{eq:curl-product}
\end{equation}

\paragraph{Divergence of a cross product.}
\begin{equation}
 \divg(\FF \times \GG) = \GG\cdot\curl\FF - \FF\cdot\curl\GG.
 \label{eq:div-cross}
\end{equation}
This identity is the vector-calculus engine behind every
integration-by-parts argument in this book: integrating both sides over
a volume and applying the divergence theorem converts a volume integral
of curls into a surface integral.


%% ================================================================
\section{Integration by parts for curl operators}
\label{sec:integration-by-parts}

Let $\FF$ and $\GG$ be smooth vector fields on a domain $\Omega$ with
boundary $\partial\Omega$ and outward normal $\nn$. Then
\begin{equation}
 \int_\Omega \GG\cdot\curl\FF\;\dd V
 = \int_\Omega \FF\cdot\curl\GG\;\dd V
 + \oint_{\partial\Omega}(\FF\times\GG)\cdot\nn\;\dd S.
 \label{eq:curl-ibp-app}
\end{equation}
This is the fundamental self-adjointness identity for the curl operator.
Setting $\FF = \EE$ and $\GG = \mur^{-1}\curl\EE$ (or the
corresponding magnetic formulation) and requiring the surface integral
to vanish yields the Hermiticity condition for the Maxwell curl--curl
operator (Chapter~\ref{ch:energy-inner}).

\subsection{Boundary conditions that kill the surface term}
\label{ssec:vanishing-surface}

The surface integral in~\eqref{eq:curl-ibp-app} vanishes if any of the
following conditions hold on $\partial\Omega$:
\begin{enumerate}
 \item \textbf{Perfect electric conductor (PEC):}
 $\nn\times\EE = 0$ on $\partial\Omega$.
 \item \textbf{Perfect magnetic conductor (PMC):}
 $\nn\times\HH = 0$ on $\partial\Omega$.
 \item \textbf{Periodic boundary conditions:}
 $\FF$ and $\GG$ satisfy Bloch periodicity on opposite faces, so the
 contributions cancel pairwise.
 \item \textbf{Fields decay sufficiently fast:}
 for an unbounded domain, if $|\FF|$ and $|\GG|$ decay faster than
 $1/r$ as $r \to \infty$, the surface integral at infinity vanishes.
\end{enumerate}
When the surface term does \emph{not} vanish---as for outgoing radiation
boundary conditions (Silver--M\"uller condition) or for modes with
$|\FF| \sim e^{+\Im k\,r}$---the curl--curl operator is no longer
self-adjoint, and the system is non-Hermitian
(Chapter~\ref{ch:boundaries}).


%% ================================================================
\section{The divergence theorem and Poynting's theorem}
\label{sec:divergence-poynting}

The divergence theorem relates a volume integral of a divergence to a
surface integral:
\begin{equation}
 \int_\Omega \divg\FF\;\dd V
 = \oint_{\partial\Omega} \FF\cdot\nn\;\dd S.
 \label{eq:divergence-theorem}
\end{equation}
Applying this to the Poynting vector
$\SS = \EE\times\HH$ gives the integral form of Poynting's theorem:
\begin{equation}
 -\pdv{}{t}\int_\Omega u\;\dd V
 = \oint_{\partial\Omega} \SS\cdot\nn\;\dd S
 + \int_\Omega \sigma\,|\EE|^2\;\dd V,
 \label{eq:poynting-integral}
\end{equation}
where $u$ is the electromagnetic energy density and $\sigma$ is the
conductivity (representing ohmic loss).


%% ================================================================
\section{Bloch-periodic boundary terms}
\label{sec:bloch-boundary}

For a photonic crystal with lattice vectors $\aa_1, \aa_2, \aa_3$ and
Bloch wave vector $\kk$, the fields satisfy
\begin{equation}
 \FF(\rr + \aa_j) = e^{i\kk\cdot\aa_j}\,\FF(\rr).
 \label{eq:bloch-bc}
\end{equation}
On a unit cell $\Omega_{\mathrm{cell}}$ with faces $S_j^+$ and $S_j^-$
related by $S_j^+ = S_j^- + \aa_j$, the surface integral
in~\eqref{eq:curl-ibp-app} cancels pairwise when the two fields belong
to compatible Bloch sectors. For example, if both fields carry the same
real Bloch wave vector and the inner product uses the usual conjugation,
or if an unconjugated bilinear form pairs $\kk$ with $-\kk$, the two
opposite faces contribute
\begin{equation}
 \int_{S_j^+}(\FF\times\GG)\cdot\nn\;\dd S
 + \int_{S_j^-}(\FF\times\GG)\cdot\nn\;\dd S = 0 .
 \label{eq:bloch-surface-cancel}
\end{equation}
The cancellation is not simply the algebraic difference
$e^{i\kk\cdot\aa_j}-e^{-i\kk\cdot\aa_j}$; it also uses the opposite
orientations of the outward normals and the adjoint boundary condition.
For real $\kk$ in a lossless periodic medium this gives the usual
self-adjoint Bloch Maxwell problem. For complex $\kk$, as in non-Bloch
or GBZ constructions, the operator is generally non-self-adjoint and
one must pair right and adjoint Bloch sectors explicitly.


%% ================================================================
\section{Silver--M\"uller radiation condition}
\label{sec:silver-muller}

For an open scatterer in three dimensions, the outgoing-wave condition
at infinity is the Silver--M\"uller condition:
\begin{equation}
 \lim_{r\to\infty} r\,\bigl[
 \hat{\rr}\times(\curl\EE) - ik\,\EE
 \bigr] = 0,
 \label{eq:silver-muller-app}
\end{equation}
uniformly in all directions $\hat{\rr}$. This condition selects
outgoing spherical waves and excludes incoming waves. It is the
mathematical statement of the ``no incoming radiation'' assumption that
makes the scattering problem well posed and the Maxwell operator
non-self-adjoint (Chapter~\ref{ch:boundaries}).


%% ================================================================
\section{Notation conventions}
\label{sec:notation}

\subsection{Time convention}
The time-harmonic convention throughout this book is
$e^{-i\omega t}$. With this choice:
\begin{itemize}
 \item A passive medium has $\Im\epsr > 0$ (absorption).
 \item A gain medium has $\Im\epsr < 0$ (amplification).
 \item A QNM eigenfrequency $\omtil = \omega_r - i\gamma$ with
 $\gamma > 0$ represents a decaying mode.
\end{itemize}

\subsection{Field quantities}
\begin{center}
\begin{tabular}{lll}
\toprule
Symbol & Name & SI Unit \\
\midrule
$\EE$ & Electric field & V/m \\
$\HH$ & Magnetic field & A/m \\
$\DD$ & Electric displacement & C/m$^2$ \\
$\BB$ & Magnetic induction & T \\
$\SS = \EE\times\HH$ & Poynting vector & W/m$^2$ \\
$\epsr$ & Relative permittivity & dimensionless \\
$\mur$ & Relative permeability & dimensionless \\
$\omega$ & Angular frequency & rad/s \\
$\omtil$ & Complex eigenfrequency & rad/s \\
$\kk$ & Wave vector & 1/m \\
$c$ & Speed of light in vacuum & m/s \\
\bottomrule
\end{tabular}
\end{center}

\subsection{Operators and inner products}
\begin{center}
\begin{tabular}{ll}
\toprule
Symbol & Meaning \\
\midrule
$\curl$ & Curl operator \\
$\divg$ & Divergence operator \\
$\grad$ & Gradient operator \\
$\hat{\Theta}$ & Maxwell curl--curl operator \\
$\Smat$ & Scattering matrix \\
$\Tmat$ & Transfer matrix \\
$\langle\cdot|\cdot\rangle$ & Standard (conjugated) inner product \\
$(\cdot,\cdot)$ & Unconjugated bilinear form \\
$\Qfac$ & Quality factor \\
$\wnum$ & Winding number \\
$\PT$ & Parity--time operator \\
\bottomrule
\end{tabular}
\end{center}


%% ================================================================
\section{Electromagnetic reciprocity from vector calculus}
\label{sec:reciprocity-vector}

The Lorentz reciprocity theorem, introduced physically in
Chapter~\ref{ch:energy-inner}, can be derived purely from the vector
calculus identities of this appendix.

\subsection{Statement and proof}
\label{ssec:reciprocity-proof}

Consider two solutions $(\EE_1, \HH_1)$ and $(\EE_2, \HH_2)$ of
Maxwell's equations at the same frequency $\omega$ in a medium with
permittivity tensor $\boldsymbol{\varepsilon}$ and permeability
$\mu_0$. The sources are localised current densities $\JJ_1$ and
$\JJ_2$. Define the \emph{reaction} of field~1 on source~2:
\begin{equation}
 \langle 1, 2 \rangle = \int_V \EE_1 \cdot \JJ_2\;\dd^3 r.
 \label{eq:reaction-app}
\end{equation}
Lorentz reciprocity states that
$\langle 1, 2 \rangle = \langle 2, 1 \rangle$ whenever
$\boldsymbol{\varepsilon} = \boldsymbol{\varepsilon}^\transpose$
(the permittivity tensor is symmetric, not necessarily Hermitian).

\paragraph{Proof.}
Start from the divergence-of-cross-product identity~\eqref{eq:div-cross}:
\begin{equation}
 \nabla \cdot (\EE_1 \times \HH_2 - \EE_2 \times \HH_1)
 = \HH_2 \cdot (\nabla \times \EE_1) - \EE_1 \cdot (\nabla \times \HH_2)
 - \HH_1 \cdot (\nabla \times \EE_2) + \EE_2 \cdot (\nabla \times \HH_1).
 \label{eq:reciprocity-divergence}
\end{equation}
Substituting Maxwell's curl equations
$\nabla \times \EE_j = i\omega\mu_0\HH_j$ and
$\nabla \times \HH_j = -i\omega\boldsymbol{\varepsilon}\EE_j + \JJ_j$
into the right-hand side:
\begin{align}
 \text{RHS} &= i\omega\mu_0\,\HH_2 \cdot \HH_1
 + i\omega\,\EE_1 \cdot \boldsymbol{\varepsilon}\EE_2
 - \EE_1 \cdot \JJ_2
 \notag\\
 &\quad - i\omega\mu_0\,\HH_1 \cdot \HH_2
 - i\omega\,\EE_2 \cdot \boldsymbol{\varepsilon}\EE_1
 + \EE_2 \cdot \JJ_1.
 \label{eq:reciprocity-rhs}
\end{align}
The $\mu_0$ terms cancel. The permittivity terms give
$i\omega(\EE_1 \cdot \boldsymbol{\varepsilon}\EE_2
- \EE_2 \cdot \boldsymbol{\varepsilon}\EE_1)$, which vanishes if and
only if $\boldsymbol{\varepsilon} = \boldsymbol{\varepsilon}^\transpose$.
What remains is
\begin{equation}
 \nabla \cdot (\EE_1 \times \HH_2 - \EE_2 \times \HH_1)
 = \EE_2 \cdot \JJ_1 - \EE_1 \cdot \JJ_2.
 \label{eq:reciprocity-simplified}
\end{equation}
Integrating over a volume $V$ that contains both sources and applying
the divergence theorem, the left-hand side becomes a surface integral
that vanishes when the fields satisfy Silver--M\"uller radiation
conditions at infinity (or when the surface is a PEC boundary):
\begin{equation}
 \oint_{\partial V} (\EE_1 \times \HH_2 - \EE_2 \times \HH_1)
 \cdot \hat{\nn}\;\dd S = 0.
 \label{eq:reciprocity-surface}
\end{equation}
Therefore $\langle 2, 1 \rangle - \langle 1, 2 \rangle = 0$, proving
Lorentz reciprocity.

\subsection{Reciprocity versus Hermiticity}
\label{ssec:reciprocity-vs-hermiticity}

The proof makes clear the distinction emphasised in
Chapter~\ref{ch:energy-inner}: reciprocity requires
$\boldsymbol{\varepsilon}^\transpose = \boldsymbol{\varepsilon}$
(symmetric), while Hermiticity of the Maxwell operator requires
$\boldsymbol{\varepsilon}^\dagger = \boldsymbol{\varepsilon}$
(Hermitian, i.e.\ real for a scalar permittivity). A medium with
complex symmetric $\boldsymbol{\varepsilon}$ (such as a lossy isotropic
dielectric) is reciprocal but not Hermitian. Conversely, a
magneto-optical medium with real antisymmetric off-diagonal elements is
Hermitian (lossless) but not reciprocal.

This distinction is central to non-Hermitian photonics: most passive
non-Hermitian systems (lossy dielectrics, absorbing metasurfaces) are
reciprocal, and reciprocity constrains their scattering matrices and
coupled-mode models. Reciprocity can fail when the material tensor,
temporal modulation, active feedback, or an effective modal reduction
lacks the relevant transpose symmetry. Such non-reciprocal or
asymmetric effective couplings are one route to point-gap winding and
the non-Hermitian skin effect of Chapter~\ref{ch:skin-effect}, although
the precise criterion is spectral rather than simply the presence or
absence of material loss.
